\chapter*{Guide du professeur}
\addcontentsline{toc}{chapter}{Guide du professeur}

\section*{1. Choix éditoriaux}

Le parcours principal réunit les notions prescrites par le PE et le RAPE T10 :
logique, algorithmique, calcul dans $\R$, équations et inéquations, fonctions de
référence, droites et cercles, trigonométrie, vecteurs et statistique descriptive.
Les contenus qui dépassent ce périmètre sont des enrichissements facultatifs :
ils portent désormais la mention explicite « non exigible T10 » et ne figurent
dans aucun sujet d'évaluation de la banque.

\begin{center}
\small
\begin{tabularx}{\textwidth}{@{}p{3.2cm}X X@{}}
\toprule
Point de vigilance & Choix du manuel & Mise en œuvre \\
\midrule
Rythme & Progression sur 33 semaines. & Deux séances par semaine, avec temps de régulation. \\
Fonctions & Priorité aux fonctions carré, cube, valeur absolue et racine carrée. &
Fonctions affine et inverse proposées en consolidation. \\
Systèmes & Réinvestissement facultatif du calcul algébrique. &
Chapitre court et clairement signalé comme extension. \\
Trigonométrie & Articuler cercle, valeurs remarquables et fonctions. &
Figures lisibles et exercices de lecture avant les calculs. \\
Statistiques & Distinguer calcul, représentation et interprétation. &
Quartiles et régression proposés pour aller plus loin. \\
\bottomrule
\end{tabularx}
\end{center}

\section*{2. Architecture d'une séquence}

La structure répétée dans l'ouvrage n'est pas un ordre rigide, mais un cycle
d'apprentissage. Pour une notion nouvelle, on recommande :

\begin{enumerate}
  \item \textbf{Réactiver} (5--10 min) : trois à cinq questions sans document sur
  les prérequis et une notion ancienne. Correction immédiate, sans note.
  \item \textbf{Faire chercher} (15--25 min) : situation courte, données simples,
  production observable. Le but est de faire émerger le besoin de la notion.
  \item \textbf{Mettre en commun} (10--15 min) : comparer deux procédures, traiter
  une erreur typique et faire préciser le vocabulaire.
  \item \textbf{Institutionnaliser} (10--20 min) : trace écrite concise contenant
  définition, propriété, conditions d'emploi et exemple minimal.
  \item \textbf{Guider} (15--25 min) : exemple complètement résolu, puis exemple à
  compléter où les aides disparaissent progressivement.
  \item \textbf{Faire pratiquer} (20--40 min) : exercices proches, puis variés,
  enfin problème de transfert. Le choix de la méthode doit progressivement
  revenir à l'élève.
  \item \textbf{Vérifier} (5 min) : billet de sortie individuel. La réponse décide
  de la reprise ou de l'étape suivante; elle ne sert pas seulement à attribuer
  une note.
\end{enumerate}

Les exercices résolus réduisent la recherche inutile au début de
l'apprentissage; les séries mixtes apprennent ensuite à reconnaître la méthode
appropriée, pas seulement à répéter le dernier algorithme vu. L'espacement et le
rappel régulier sont assurés par les automatismes de chaque séance.

\section*{3. Format d'une séance de deux heures}

\begin{center}
\small
\begin{tabularx}{\textwidth}{@{}>{\raggedright\arraybackslash}p{2.7cm}
  >{\raggedright\arraybackslash}p{2.4cm}X@{}}
\toprule
Phase & Durée indicative & Production attendue \\
\midrule
Automatismes & 10 min & Réponses individuelles, correction diagnostique. \\
Rappel ciblé & 10 min & Une définition et un exemple ancien réactivés. \\
Recherche & 25 min & Brouillon de groupe ou individuel, conjecture justifiée. \\
Mise en commun & 15 min & Comparaison de procédures et traitement d'une erreur. \\
Trace écrite & 20 min & Résultat précis, conditions d'emploi, exemple. \\
Pratique guidée & 20 min & Exemple partiellement complété puis exercice analogue. \\
Pratique autonome & 15 min & Deux exercices différenciés ou un problème court. \\
Billet de sortie & 5 min & Une question essentielle, rendue avant de quitter la salle. \\
\bottomrule
\end{tabularx}
\end{center}

Pour une séance de 55 minutes, conserver automatismes, un objectif unique,
pratique et billet de sortie; répartir recherche et institutionnalisation sur
deux séances au lieu de les comprimer.

\section*{4. Progression annuelle opérationnelle}

Chaque semaine correspond à quatre heures. Les deux colonnes peuvent représenter
deux séances de deux heures ou être redistribuées selon l'emploi du temps. Les
notions sont réintroduites après un délai afin de favoriser leur mémorisation
durable.

\begingroup
\footnotesize
\begin{longtable}{@{}c >{\raggedright\arraybackslash}p{5.1cm}
  >{\raggedright\arraybackslash}p{5.1cm}
  >{\raggedright\arraybackslash}p{2.6cm}@{}}
\toprule
Sem. & Séance A & Séance B & Vérification \\
\midrule
\endhead
1 & Énoncés, propositions, vérité & Connecteurs « et », « ou », négation & Billet logique \\
2 & Implication, équivalence & Quantificateurs et négations & Quiz 1 \\
3 & Variables, affectation, séquence & Lire et exécuter un pseudo-code & Trace d'algorithme \\
4 & Conditions & Boucles, organigrammes & Mini-tâche \\
5 & Ensembles $\N,\Z,\D,\Q,\R$ & Fractions et puissances & Automatismes 1 \\
6 & Racines carrées et ordre & Intervalles, distance, valeur absolue & Quiz 2 \\
7 & Droite : équations cartésienne et réduite & Droite : équation paramétrique & Banque : exercices ciblés \\
8 & Écritures et valeurs approchées & Encadrements et propagation & Problème court \\
9 & Forme canonique & Discriminant et nombre de racines & Automatismes 2 \\
10 & Résolution du second degré & Factorisation du trinôme & Quiz 3 \\
11 & Signe du trinôme & Inéquations du second degré & Billet de sortie \\
12 & Cubique avec racine connue & Division et identification & Exercice raisonné \\
13 & Signe et inéquations cubiques & Équations rationnelles & Quiz 4 \\
14 & Cercle : équations & Intersection droite-cercle & Banque : problème de géométrie \\
15 & Notion de fonction, domaine, image & Lecture graphique & Automatismes 3 \\
16 & Croissance, décroissance, extremum & Parité & Quiz 5 \\
17 & Fonction carré & Fonction cube & Tracé contrôlé \\
18 & Valeur absolue & Racine carrée & Banque : évaluation formative \\
19 & Degrés et radians & Cercle trigonométrique & Automatismes 4 \\
20 & Angles remarquables & Relation $\cos^2x+\sin^2x=1$ & Quiz 6 \\
21 & Parité et périodicité trigonométriques & Réinvestissement fonctions-trigonométrie & Banque : exercice structuré \\
22 & Courbes de référence : lecture & Tables et tableaux de variation & Billet graphique \\
23 & Problèmes sur les fonctions & Pratique mixte différenciée & Quiz 7 \\
24 & Vecteurs : sens, direction, norme & Égalité et opérations & Construction \\
25 & Coordonnées et milieu & Colinéarité & Automatismes 5 \\
26 & Alignement et parallélisme & Produit scalaire coordonné & Quiz 8 \\
27 & Orthogonalité & Problèmes vectoriels & Banque : problème de synthèse \\
28 & Population, variable discrète/continue & Tableaux d'effectifs et fréquences & Collecte réelle \\
29 & Moyenne et médiane & Variance et écart-type & Automatismes 6 \\
30 & Diagramme en bâtons et circulaire & Histogramme & Production graphique \\
31 & Interpréter et critiquer une représentation & Problème statistique complet & Banque : sommative \\
32 & Révisions mixtes & Composition libre depuis la banque & Bilan individuel \\
33 & Remédiation ciblée & Extensions/olympiades/projet & Plan de progrès \\
\bottomrule
\end{longtable}
\endgroup

\section*{5. Différenciation sans abaisser les objectifs}

\begin{itemize}
  \item \textbf{Soutien :} réduire le nombre de données, fournir une figure ou une
  première étape, laisser visibles les conditions d'emploi d'une méthode.
  \item \textbf{Parcours standard :} exercices \diff{1}, puis \diff{2}, avec une
  question de justification explicite.
  \item \textbf{Approfondissement :} retirer les indications, modifier une donnée,
  demander une généralisation ou proposer un défi \diff{3}.
  \item \textbf{Élèves très avancés :} privilégier les preuves, invariants,
  contre-exemples et problèmes olympiques; ne pas seulement augmenter la longueur
  des calculs.
\end{itemize}

Le travail en groupe n'est utile que si la responsabilité individuelle reste
visible : chacun prépare d'abord une réponse, un rapporteur est tiré au sort, et
un billet de sortie individuel clôt la séance.

\section*{6. Évaluation et retour d'information}

La \emph{Banque d'exercices et de problèmes} permet de composer des évaluations
adaptées à la séquence effectivement enseignée. Chaque bloc est autonome,
structuré en sous-questions et accompagné d'une indication relative à la
calculatrice.

\begin{center}
\small
\begin{tabularx}{\textwidth}{@{}p{2.5cm}p{5.2cm}X@{}}
\toprule
Usage & Composition conseillée & Retour attendu \\
\midrule
Diagnostique & Exercices de domaines différents portant sur les prérequis. & Repérer les notions à reprendre. \\
Formative & Exercices ciblés, indépendants et limités aux capacités travaillées. & Correction rapide et reprise ciblée. \\
Sommative & Exercices de plusieurs domaines et problème de synthèse. & Bilan d'une période et critères explicites. \\
\bottomrule
\end{tabularx}
\end{center}

Les formatives servent d'abord à réguler l'enseignement; les sommatives attestent
les acquis d'une période. Les structures BEPC, Bac et internationales fournissent
des repères d'organisation : questions courtes, exercices progressifs et problème
de synthèse.

Les contenus restent strictement ceux de T10. Les apports internationaux retenus
sont l'indépendance des exercices, une progression des réponses brèves vers les
tâches étendues, l'affichage des points et la valorisation de la justification,
de la modélisation et de la communication.


Une méthode correcte reste valorisée même si une erreur de calcul ultérieure
affecte le résultat. Inversement, un résultat exact sans justification ne reçoit
pas les points réservés au raisonnement.

\section*{7. Références pédagogiques utilisées}

\begin{itemize}
  \item Education Endowment Foundation, \emph{Mathematics in Key Stages 2 and 3:
  Evidence Review}, 2018. Revue de plus de 3 000 études :
  \href{https://educationendowmentfoundation.org.uk/education-evidence/evidence-reviews/mathematics-in-key-stages-2-and-3}{revue disponible en ligne}.
  \item H. Pashler et al., \emph{Organizing Instruction and Study to Improve
  Student Learning}, U.S. Institute of Education Sciences, 2007 :
  \url{https://ies.ed.gov/ncee/wwc/practiceguide/1}.
  \item J. Sweller et G. A. Cooper, « The Use of Worked Examples as a Substitute
  for Problem Solving in Learning Algebra », \emph{Cognition and Instruction},
  2(1), 1985, p. 59--89, DOI: \url{https://doi.org/10.1207/s1532690xci0201_3}.
  \item D. Rohrer et K. Taylor, « The Shuffling of Mathematics Problems Improves
  Learning », \emph{Instructional Science}, 35, 2007, p. 481--498,
  DOI: \url{https://doi.org/10.1007/s11251-007-9015-8}.
  \item P. Black et D. Wiliam, « Assessment and Classroom Learning »,
  \emph{Assessment in Education}, 5(1), 1998, p. 7--74,
  DOI: \url{https://doi.org/10.1080/0969595980050102}.
\end{itemize}

Ces références justifient des choix de structure; elles ne remplacent ni le PE
ni le jugement professionnel de l'enseignant dans son contexte de classe.
